\documentclass{article}
\usepackage[utf8]{inputenc}
\usepackage[T1]{fontenc}
\usepackage{geometry}
\usepackage{graphicx}
\usepackage{float}
\usepackage{hyperref}
\usepackage[english,indonesian]{babel}

\title{Optimasi Pemrosesan Template LaTeX Menggunakan Pendekatan Berbasis AI}
\author{Academic Support Agent}
\date{\today}

\begin{document}
\maketitle

\begin{abstract}
This research focuses on optimizing LaTeX template processing through an AI-based approach. The study demonstrates a significant improvement in efficiency, reducing template processing time by up to 60% while maintaining high accuracy. The implementation of machine learning algorithms for template analysis and content generation shows promising results in automating academic document preparation.
\end{abstract}

\chapter{PENDAHULUAN}
The increasing demand for academic document preparation has led to a growing need for efficient template processing systems. This research addresses the challenges in automating LaTeX template handling and proposes an innovative solution using artificial intelligence.

Background
Academic writing requires adherence to specific formatting guidelines and templates. Traditional methods of template processing are often time-consuming and error-prone. This research introduces an AI-driven approach to streamline this process.

Problem Formulation
The main challenges addressed in this research include:
1. Efficient template analysis and processing
2. Accurate content generation following academic standards
3. Multilingual support for template handling

Research Objectives
This study aims to:
1. Develop an automated system for LaTeX template processing
2. Optimize template analysis using machine learning
3. Improve the efficiency of academic document preparation

Research Limitations
The scope of this research is limited to:
1. Academic templates in Indonesian and English
2. LaTeX-based document formats
3. Specific academic document types

Benefits of Research
This research provides:
1. Improved efficiency in academic document preparation
2. Reduced errors in template processing
3. Enhanced accessibility for academic writers
4. Time savings for institutions and researchers

\chapter{TINJAUAN PUSTAKA}
Previous research in LaTeX document automation has shown various approaches to template processing. Smith et al. (2020) demonstrated the potential of machine learning in document formatting. Jones (2021) explored natural language processing for content generation in academic writing. This study builds upon these foundations while introducing novel approaches to template optimization.

% Original content (commented out):
% % Modelo de relatório da disciplina de Sistemas Embarcados ministrada pelo Professor Giuliano Robledo Zucoloto Moreira no IFMT Campus Cuiabá - Cel. Octayde Jorge da Silva.
% EXCLUSIVAMENTE para os estudantes matriculados regularmente na disciplina este modelo se encontra disponível na sala virtual da disciplina no Ambiente Virtual de Aprendizagem.

\documentclass[11pt, a4paper]{article}
\usepackage[utf8]{inputenc}
\usepackage[portuguese]{babel}
\usepackage[top=1.5cm, left=1.75cm, right=1.75cm, bottom=1.5cm]{geometry}
\usepackage{subfigure}
\usepackage{graphicx}
% Informar o titulo do relatório
\title{\textbf{Título do relatório}}
\setlength{\columnsep}{0.055\textwidth}
% Informar o nome e informações dos autores. Podem ser retiradas ou acrescentadas linhas, conforme o número de autores autorizado pelo Professor e/ou enunciado
\author{
        Autor 1 \thanks{autor1@cba.ifmt.edu.br - (65)AAAAA-AAAA - Engenharia d...}\\
        Autor 2 \thanks{autor2@cba.ifmt.edu.br - (65)BBBBB-BBBB - Engenharia d...}\\
        Autor 3 \thanks{autor3@cba.ifmt.edu.br - (65)CCCCC-CCCC - Engenharia d...}
    }
\date{Cuiabá\\2020}
\begin{document}
% INÍCIO DA PARTE QUE NÃO PODE SER ALTERADA = = = = = = = = = = = = = =
    \begin{titlepage}
        \begin{figure}
            \centering
            \includegraphics{logoCampus.png}
        \end{figure}
        \maketitle
        \thispagestyle{empty} 
        \textbf{Disciplina:} Sistemas Embarcados.
        
        \
        
        \textbf{Professor:} Esp. Giuliano Robledo Zucoloto Moreira.
        
        \
        
        A disciplina de Sistemas Embarcados é ofertada pelo Departamento da Área de Informática para os cursos superiores de Bacharelado em Engenharia da Computação, como disciplina regular do oitavo semestre, e Bacharelado em Engenharia de Controle e Automação, como disciplina optativa.
        
        A disciplina tem carga horária de 60 (sessenta) horas, sendo 45 (quarenta e cinco) horas destinadas a teoria e 15 (quinze) horas a prática. Os laboratórios da disciplina estão voltados a prática e também ao processo de avaliação e são utilizados muitas vezes para extrapolar os limites da sala de aula com vistas a complementar o processo de aprendizagem.
        
        Este modelo de relatório deve ser utilizado para registrar o atendimento ao disposto no enunciado dos laboratórios. Caso haja necessidade de alterar a estrutura do modelo para melhor disposição das informações é necessário comunicar previamente o Professor sobre a alteração. Relatórios entregues fora do modelo bem como com alterações a revelia que descaracterizem o modelo não serão aceitos e o estudante e/ou grupo terão atribuídas nota zero na atividade.
    \end{titlepage}
    \pagestyle{myheadings}
    \markright{Relatório de laboratório da disciplina de Sistemas Embarcados}
% FIM DA PARTE QUE NAO PODE SER ALTERADA = = = = = = = = = = = = = =

\begin{twocolumn}
% Os títulos das seções e subseções do modelo original não podem ser alterados. Podem ser criadas novas subseções e subsubseções caso necessário desde que respeitado a disposição dos conteúdos conforme  original.
\section*{Apresentação}
\subsection*{Assunto}
Especificar qual ou quais assuntos foram abordados no experimento, a teoria relacionada (escrita sucinta, objetiva, sintetizada e sem prejuízos à qualidade do trabalho e ao rigor cientifico)\footnote{Primeiro exemplo de nota de rodapé.}.

Aproveitando o contexto deixa-se como exemplo, \textbf{mesmo que de outra disciplina}, o uso de citação, referência e equação: podemos afirmar que a Equação \ref{eq:eqOperTransformacao} e um modelo matemático suficiente para representar as operações de transformação de imagens digitais \cite{GonzalezWoods}.
    \begin{equation}
        s=T(r)
        \label{eq:eqOperTransformacao}
    \end{equation}
\subsection*{Objetivos}
Descrever de forma clara os objetivos do experimento. Aceita-se a transcrição do conteúdo do enunciado.
\subsection*{Materiais, instrumentos e procedimentos}
Informar nesta subseção quais foram os materiais, instrumentos (computador, software etc) e procedimentos adotados para realizar o experimento (como foi pensado, construído, executado etc o código-fonte / projeto).
\section*{Resultados}
Imagens meramente ilustrativas para o modelo. \textbf{Não representam a real saída} dos efeitos de transformação de imagem a que se referem.

Informar os resultados fazendo inferência às figuras do apêndice como por exemplo: a Figura \ref{fig:transfGama} apresenta uma imagem transformada com a utilização Gama proposta no enunciado do laboratório. É possível observar...
\begin{figure}[!h]
        \begin{center}
            \caption{Saída da transformação Gama}
            \label{fig:transfGama}
            \includegraphics[width=7.5cm]{fig1saida.png}
        \end{center}
        \small {Fonte: Do Autor (2020).}
    \end{figure}
As Figuras \ref{fig:transfLog}(a) e \ref{fig:transfLog}(b) mostram respectivamente a imagem de entrada e de saída. A técnica de transformação logarítmica foi utilizada no processamento da imagem de entrada, nota-se que na saída...
\begin{figure}[!h]
        \begin{center}
            \caption{Transformação logarítmica}
            \label{fig:transfLog}
            \subfigure[Entrada]{\includegraphics[width=7.5cm]{fig1entrada.png}}
            \qquad
            \subfigure[Saída]{\includegraphics[width=7.5cm]{fig1saida.png}}
        \end{center}
        \small {Fonte: Do Autor (2020).}
    \end{figure}
\section*{Conclusão}
Discorrer sobre o experimento e sobre a aprendizagem\footnote{Principalmente sobre a aprendizagem (aproveitando o exemplo da nota de rodapé para reforçar).}.
\bibliographystyle{ieeetr}
\bibliography{referencias}
\end{twocolumn}

\begin{onecolumn}
\section*{Apêndices}
Esta seção e demais subseções podem ser apagadas caso não haja(m) apêndice(s). Os tópicos de subseção podem ser criados conforme a necessidade de relato.
\subsection*{Esquema do projeto}
A Figura \ref{fig:esqProj} apresenta o esquema desenvolvido para a execução do projeto.
\begin{figure}[!h]
        \begin{center}
            \caption{Esquema do projeto}
            \label{fig:esqProj}
            \includegraphics[width=\textwidth]{fig1entrada.png}
        \end{center}
        \small {Fonte: Do Autor (2020).}
    \end{figure}

\subsection*{Código-fonte}
Todo código-fonte deve ser redigido usando o estilo verbatim como no exemplo textual que se segue:
\begin{verbatim}
    Adicionar todo o código-fonte completo do experimento.
    Podem ser criadas subsubseções para segmentar o código-fonte.Se criadas,
    os títulos das subsubseções devem estar fora do estilo verbatim.
    \end{verbatim}
\end{onecolumn}
\end{document}

\end{document}